% Clase 8 --- Por qué conviene el potencial: es un ESCALAR (§3, §4).
% Para el campo hay que descomponer cada aporte y sumar componente a
% componente; para el potencial se suman números. Y del potencial se recupera
% el campo (próxima clase), así que no se pierde información.
\begin{center}
\begin{tikzpicture}[scale=1]

  % =============== (a) campo: suma vectorial ===============
  \begin{scope}
    \node[figpos] at (-1.5,1.25) {};
    \node[figlbl, anchor=east] at (-1.65,1.25) {$Q_1$};
    \node[figpos] at (-1.7,-0.9) {};
    \node[figlbl, anchor=east] at (-1.85,-0.9) {$Q_2$};
    \node[figneg] at (0.65,-1.5) {};
    \node[figlbl, anchor=north] at (0.65,-1.7) {$Q_3$};

    \node[figneutra, minimum size=5pt] (P) at (1.15,0.55) {};
    \node[figlbl, anchor=east] at (1.05,0.62) {$P$};

    % aportes vectoriales
    \draw[figvec] (P) -- ++(40:1.15);
    \draw[figvec] (P) -- ++(28:0.85);
    \draw[figvec] (P) -- ++(-105:0.95);
    \node[figlblaux, anchor=south west] at (2.05,1.35) {$\vec E_1$};
    \node[figlblaux, anchor=north west] at (1.95,1.0) {$\vec E_2$};
    \node[figlblaux, anchor=west] at (1.3,-0.45) {$\vec E_3$};

    \node[figlbl, align=center] at (0,-2.85)
      {$\vec E = \sum_i \vec E_i$: \ hay que \textbf{proyectar}\\
       cada aporte y sumar componentes};
    \node[figlbl] at (0,-3.95) {(a) campo eléctrico};
  \end{scope}

  % =============== (b) potencial: suma de números ===============
  \begin{scope}[xshift=7.6cm]
    \node[figpos] at (-1.5,1.25) {};
    \node[figlbl, anchor=east] at (-1.65,1.25) {$Q_1$};
    \node[figpos] at (-1.7,-0.9) {};
    \node[figlbl, anchor=east] at (-1.85,-0.9) {$Q_2$};
    \node[figneg] at (0.65,-1.5) {};
    \node[figlbl, anchor=north] at (0.65,-1.7) {$Q_3$};

    \node[figneutra, minimum size=5pt] (P2) at (1.15,0.55) {};
    \node[figlbl, anchor=south] at (1.15,0.75) {$P$};

    % sólo distancias
    \draw[figvecaux] (-1.5,1.25) -- (P2);
    \draw[figvecaux] (-1.7,-0.9) -- (P2);
    \draw[figvecaux] (0.65,-1.5) -- (P2);
    \node[figlblaux, anchor=south] at (-0.2,0.98) {$R_1$};
    \node[figlblaux, anchor=south east] at (-0.2,-0.28) {$R_2$};
    \node[figlblaux, anchor=west] at (1.0,-0.5) {$R_3$};

    \node[figlbl, align=center] at (0,-2.85)
      {$\displaystyle V = \frac{1}{4\pi\varepsilon_0}\sum_i \frac{Q_i}{R_i}$:
       \ se suman \textbf{números}};
    \node[figlbl] at (0,-3.95) {(b) potencial};
  \end{scope}

\end{tikzpicture}\\[0.5em]
{\small\itshape Los dos paneles tienen \textbf{la misma} configuración de
cargas. Para el campo hacen falta direcciones, ángulos y proyecciones; para el
potencial alcanza con las distancias y el signo de cada carga. Y no se pierde
nada: el campo se recupera derivando el potencial (próxima clase). Por eso
conviene calcular primero $V$ y después $\vec E$, y no al revés.}
\end{center}
